\documentclass[aps,preprint]{revtex4}%
\usepackage{amssymb}
\usepackage{amsfonts}
\usepackage{amsmath}
\usepackage{graphicx}%
\setcounter{MaxMatrixCols}{30}
%TCIDATA{OutputFilter=latex2.dll}
%TCIDATA{Version=5.50.0.2960}
%TCIDATA{CSTFile=revtex4.cst}
%TCIDATA{Created=Thursday, April 21, 2022 09:23:27}
%TCIDATA{LastRevised=Thursday, June 09, 2022 12:19:15}
%TCIDATA{<META NAME="GraphicsSave" CONTENT="32">}
%TCIDATA{<META NAME="SaveForMode" CONTENT="1">}
%TCIDATA{BibliographyScheme=Manual}
%TCIDATA{<META NAME="DocumentShell" CONTENT="Articles\SW\REVTeX 4">}
%BeginMSIPreambleData
\providecommand{\U}[1]{\protect\rule{.1in}{.1in}}
%EndMSIPreambleData
\newtheorem{theorem}{Theorem}
\newtheorem{acknowledgement}[theorem]{Acknowledgement}
\newtheorem{algorithm}[theorem]{Algorithm}
\newtheorem{axiom}[theorem]{Axiom}
\newtheorem{claim}[theorem]{Claim}
\newtheorem{conclusion}[theorem]{Conclusion}
\newtheorem{condition}[theorem]{Condition}
\newtheorem{conjecture}[theorem]{Conjecture}
\newtheorem{corollary}[theorem]{Corollary}
\newtheorem{criterion}[theorem]{Criterion}
\newtheorem{definition}[theorem]{Definition}
\newtheorem{example}[theorem]{Example}
\newtheorem{exercise}[theorem]{Exercise}
\newtheorem{lemma}[theorem]{Lemma}
\newtheorem{notation}[theorem]{Notation}
\newtheorem{problem}[theorem]{Problem}
\newtheorem{proposition}[theorem]{Proposition}
\newtheorem{remark}[theorem]{Remark}
\newtheorem{solution}[theorem]{Solution}
\newtheorem{summary}[theorem]{Summary}
\newenvironment{proof}[1][Proof]{\noindent\textbf{#1.} }{\ \rule{0.5em}{0.5em}}
\begin{document}
\preprint{HEP/123-qed}
\title[Short title for running header]{REV\TeX4}
\author{First Author}
\affiliation{My Institution}
\author{Second Author}
\affiliation{My Institution}
\author{Third Author}
\affiliation{Other Institution}
\keywords{one two three}
\pacs{PACS number}

\begin{abstract}
Shell document for REV\TeX{} 4.

\end{abstract}
\volumeyear{year}
\volumenumber{number}
\issuenumber{number}
\eid{identifier}
\date[Date text]{date}
\received[Received text]{date}

\revised[Revised text]{date}

\accepted[Accepted text]{date}

\published[Published text]{date}

\startpage{101}
\endpage{102}
\maketitle
\tableofcontents


\section{The Pfaffian of Skew Linear Transformations}

\subsubsection{The Pfaffian}

Let $V$ be an inner product space of dimension $n=2m$ and let $G_{1}%
,\cdots,G_{n}$ $\in\mathcal{B}\left(  V\right)  $ be skew linear
transformations. Every $G_{k}$ determines an element $g_{k}=\Omega\left(
G_{k}\right)  \in%
%TCIMACRO{\tbigwedge ^{2}}%
%BeginExpansion
{\textstyle\bigwedge^{2}}
%EndExpansion
V$ by
\begin{gather}
\Omega:\mathcal{B}_{as}\left(  V\right)  \rightarrow%
%TCIMACRO{\tbigwedge \nolimits^{2}}%
%BeginExpansion
{\textstyle\bigwedge\nolimits^{2}}
%EndExpansion
V\\
\Omega\left(  \sum_{1\leq j,k\leq2m}\left\vert e_{j}\right\rangle \left\langle
e_{k}\right\vert G_{jk}\right)  =\sum_{1\leq j,k\leq2m}\left\vert e_{j}\wedge
e_{k}\right\rangle g_{jk}.
\end{gather}
Thus,
\begin{equation}
g_{1}\wedge\cdots\wedge g_{n}\in%
%TCIMACRO{\tbigwedge \nolimits^{n}}%
%BeginExpansion
{\textstyle\bigwedge\nolimits^{n}}
%EndExpansion
V.
\end{equation}


\begin{definition}
Now fix a basis vector $a$ of $%
%TCIMACRO{\tbigwedge \nolimits^{n}}%
%BeginExpansion
{\textstyle\bigwedge\nolimits^{n}}
%EndExpansion
V$ and set
\begin{equation}
g_{1}\wedge\cdots\wedge g_{n}=\Omega\left(  G_{1}\right)  \wedge\cdots
\wedge\Omega\left(  G_{n}\right)  =\frac{1}{\left\Vert a\right\Vert ^{2}%
}\operatorname{Pf}\left(  \left\{  G_{1},\cdots,G_{n}\right\}  \right)
a\equiv\frac{1}{\left\Vert a\right\Vert ^{2}}\operatorname{Pf}\left(  \left\{
g_{1},\cdots,g_{n}\right\}  \right)  a.
\end{equation}
The scalar $\operatorname{Pf}\left(  \left\{  G_{1},\cdots,G_{n}\right\}
\right)  $ is called the Pfaffian of $\left\{  G_{1},\cdots,G_{n}\right\}  $.
Since every two elements $\left\{  g_{j},g_{k}\right\}  $ commute, it follows
that $\operatorname{Pf}$ is a \emph{symmetric} $m$-linear function in
\begin{equation}
\operatorname{Pf}:\overset{m}{\overbrace{%
%TCIMACRO{\tbigwedge \nolimits^{2}}%
%BeginExpansion
{\textstyle\bigwedge\nolimits^{2}}
%EndExpansion
V\times\cdots\times%
%TCIMACRO{\tbigwedge \nolimits^{2}}%
%BeginExpansion
{\textstyle\bigwedge\nolimits^{2}}
%EndExpansion
V}}\rightarrow\mathbb{C}%
\end{equation}

\end{definition}

\begin{definition}
The Pfaffian of a single skew transformation $G$ is
\begin{equation}
\operatorname{Pf}\left(  G\right)  =\frac{1}{m!}\operatorname{Pf}\left(
G,\cdots,G\right)
\end{equation}
and can expressed in the form
\begin{equation}
\operatorname{Pf}\left(  G\right)  =\left\langle \left.  \Omega\left(
G\right)  ^{m}\right\vert a\right\rangle .
\end{equation}

\end{definition}

\begin{proposition}
1:If $\mathfrak{\tau}$ is an isometry of $V$, then
\begin{equation}
\operatorname{Pf}\left(  \left\{  \mathfrak{\tau}G_{1}\mathfrak{\tau}%
^{-1},\cdots,\mathfrak{\tau}G_{n}\mathfrak{\tau}^{-1}\right\}  \right)
=\det\left(  \mathfrak{\tau}\right)  \operatorname{Pf}\left(  \left\{
G_{1},\cdots,G_{n}\right\}  \right)
\end{equation}

\end{proposition}

\begin{proposition}
2:Let $a$ and $b$ be basis vectors of $V$. Then%
\begin{equation}
\operatorname{Pf}_{a}\left(  G\right)  \operatorname{Pf}_{b}\left(  G\right)
=\left\langle \left.  a\right\vert b\right\rangle \det\left(  G\right)  .
\end{equation}

\end{proposition}

\begin{corollary}%
\begin{equation}
\operatorname{Pf}_{a}\left(  G\right)  ^{2}=\left\Vert a\right\Vert ^{2}%
\det\left(  G\right)  .
\end{equation}

\end{corollary}

\subsubsection{Direct Sums}

Let $E$ and $F$ be inner product spaces. Then an inner product is induced in
the direct sum $E\oplus F$ by%
\begin{equation}
\left\langle \left.  x_{1}\oplus y_{1}\right\vert x_{2}\oplus y_{2}%
\right\rangle =\left\langle \left.  x_{1}\right\vert y_{1}\right\rangle
+\left\langle \left.  x_{2}\right\vert y_{2}\right\rangle
\end{equation}


Let%

\begin{align}
i_{E}  &  :E\rightarrow E\oplus F\\
i_{F}  &  :F\rightarrow E\oplus F
\end{align}
and
\begin{align}
\pi_{E}  &  :E\oplus F\rightarrow E\\
\pi_{F}  &  :E\oplus F\rightarrow F
\end{align}
denote the natural inclusions and projections.

\begin{proposition}
3:Let $E$ and $F$ be inner product spaces of dimensions $2p$ and $2q$
respectively. hoose basis vectors $a$ and $b$ for ${\LARGE \wedge}^{2p}E$ and
${\LARGE \wedge}^{2q}F$ then
\begin{equation}
\operatorname{Pf}_{a\oplus b}\left(  G_{E}\oplus G_{F}\right)
=\operatorname{Pf}_{a}\left(  G_{E}\right)  \operatorname{Pf}_{b}\left(
G_{F}\right)  ,\quad G_{E}\in\mathcal{B}_{as}\left(  E\right)  ,G_{F}%
\in\mathcal{B}_{as}\left(  F\right)  .
\end{equation}

\end{proposition}

\subsubsection{ Euclidean Spaces}

\begin{definition}
A Euclidean vector space is a finite-dimensional inner product space over the
real numbers.
\end{definition}

\begin{definition}
An Affine space is a set $\left\{  A\right\}  $ together with a vector space
$A$ and a transitive and free action of the additive group of $A_{+}$ on the
set $\left\{  A\right\}  $. The elements of the affine space $A$ are called
points. The vector space $A$ is said to be associated to the affine space, and
its elements are called vectors, translations, or sometimes free vectors.
\end{definition}

\begin{definition}
Let $V$ be a finite-dimensional real vector space and let $\mathbf{b}_{1}$ and
$\mathbf{b}_{2}$ be two ordered bases for $V$. It is a standard result in
linear algebra that there exists a unique linear transformation
$A:V\rightarrow V$ that takes $\mathbf{b}_{1}$ to $\mathbf{b}_{2}$. The bases
$\mathbf{b}_{1}$ and $\mathbf{b}_{2}$ are said to have the same orientation
(or be consistently oriented) if $A$ has a positive determinant; otherwise
they have opposite orientations. The property of having the same orientation
defines an equivalence relation on the set of all ordered bases for $V$. If
$V$ is non-zero, there are precisely two equivalence classes determined by
this relation. An orientation on $V$ is an assignment of $+1$ to one
equivalence class and $-1$ to the other.
\end{definition}

Let $E$ be an oriented Euclidean space of dimension $n=2m$ and let $e$ be the
unique unit vector in ${\LARGE \wedge}^{2m}E$ which represents the
orientation. Then we set%

\begin{equation}
\operatorname{Pf}_{e}\left(  \left\{  G_{1},\cdots,G_{n}\right\}  \right)
=\left\langle \left.  \Omega\left(  G_{1}\right)  \wedge\cdots\wedge
\Omega\left(  G_{m}\right)  \right\vert e\right\rangle ;\quad G_{j}%
\in\mathcal{B}_{as}\left(  E\right)  ,1\leq j\leq m
\end{equation}
and
\begin{equation}
\operatorname{Pf}_{e}\left(  G\right)  =\frac{1}{m!}\operatorname{Pf}%
_{e}\left(  \left\{  \overset{m}{\overbrace{G,\cdots,G}}\right\}  \right)
=\left\langle \left.  \Omega\left(  G\right)  \wedge\cdots\wedge\Omega\left(
G\right)  \right\vert e\right\rangle ;\quad G\in\mathcal{B}_{as}\left(
E\right)  .
\end{equation}


Observe that if the orientation of $E$ is reversed then $\operatorname{Pf}%
_{e}$ is changed to $-\operatorname{Pf}_{e}$ .

\begin{remark}
Proposition 1 shows that the Pfaffian is invariant under proper rotations.
\end{remark}

\begin{remark}
Proposition 2 implies that%
\begin{equation}
\operatorname{Pf}_{e}\left(  G\right)  ^{2}=\det\left(  G\right)  ;\quad
G\in\mathcal{B}_{as}\left(  E\right)  .
\end{equation}

\end{remark}

\begin{remark}
Finally, by Proposition 3%
\begin{equation}
\operatorname{Pf}_{e\oplus e}\left(  G_{1}\oplus G_{2}\right)
=\operatorname{Pf}_{e}\left(  G_{1}\right)  \operatorname{Pf}_{e}\left(
G_{2}\right)  ,\quad G_{1},G_{1}\in\mathcal{B}_{as}\left(  E\right)  .
\end{equation}

\end{remark}

\begin{remark}
Let $V$ be an $n$-dimensional complex vector space with a Hermitian inner
product $\left\langle \left.  {}\right\vert \right\rangle $ and let
$V_{\mathbb{R}}$ denote the underlying real vector space. Give $V_{\mathbb{R}%
}$ the natural orientation and the positive definite inner product
\end{remark}

\subsubsection{%
\begin{equation}
(\left.  x\right\vert y)=\operatorname{Re}\left(  \left\langle \left.
a\right\vert b\right\rangle \right)  ,\,x,y\in E.
\end{equation}
Canonical Form}

\begin{enumerate}
\item
\begin{equation}
g_{k}=\sum_{1\leq j\leq r}\lambda_{kj}\left\vert e_{kj}\wedge e_{kj+r}%
\right\rangle =\sqrt{2}\sum_{1\leq j\leq r}\lambda_{kj}\left\vert
e_{kj}e_{kj+r}\right\rangle ,\,\lambda_{j}\in\mathbb{R}_{+},1\leq j\leq r
\end{equation}


\item
\begin{equation}
\mathcal{R}_{e}\left(  \Omega^{-1}\left(  g_{k}\right)  \right)  =\left(
\begin{array}
[c]{cc}%
\mathbb{O}_{r} & \mathbf{\lambda}_{k}\\
-\mathbf{\lambda}_{k} & \mathbb{O}_{r}%
\end{array}
\right)  =\mathcal{R}_{e}\left(  G_{k}\right)
\end{equation}


\item
\begin{equation}
g_{1}\wedge\cdots\wedge g_{n}=\Omega\left(  G_{1}\right)  \wedge\cdots
\wedge\Omega\left(  G_{n}\right)  =\frac{1}{\left\Vert a\right\Vert ^{2}%
}\operatorname{Pf}\left(  \left\{  G_{1},\cdots,G_{n}\right\}  \right)
a\equiv\frac{1}{\left\Vert a\right\Vert ^{2}}\operatorname{Pf}\left(  \left\{
g_{1},\cdots,g_{n}\right\}  \right)  a.
\end{equation}

\end{enumerate}


\end{document}